\section*{第 \ref{ch:b3:milky-way-galaxies} 章——银河系与星系}

\begin{solution}{exo:b3:milky-way-galaxies:1}
(a)~$2\pi R_0 = 2\pi\times\qty{2.523e17}{km} = \qty{1.585e18}{km}$；以
\qty{229}{km.s^{-1}} 的速度，需时 \qty{6.92e15}{s}，即
\qty{219}{Ma}。(b)~$4540/219 = 20.7$：约二十圈。
\end{solution}

\begin{solution}{exo:b3:milky-way-galaxies:2}
(a)~$\qty{1}{au}/\qty{1}{\text{年}} = \num{1.496e8}/\num{3.156e7} =
\qty{4.74}{km.s^{-1}}$. (b)~$v = 4.74\times6.411\times8.178 =
\qty{248.5}{km.s^{-1}}$，符合综述的
\qty[separate-uncertainty=true]{248 \pm 3}{km.s^{-1}}. (c)~$248.5 - 11 = \qty{237.5}{km.s^{-1}}$,
比红巨星的 \qty{229}{km.s^{-1}} 约高 \qty{8}{km.s^{-1}}：
两者所依据的 $R_0$ 取值不同（8.178 与 \qty{8.122}{kpc}），
而\omterm{def:b3:milky-way-galaxies:lsr}{本地静止标准}本身，相对于
真正的圆轨道，可有高达约 \qty{15}{km.s^{-1}} 的漂移。
\end{solution}

\begin{solution}{exo:b3:milky-way-galaxies:3}
(a)~$n = 10(1 - 0.6) = 4$：E4。(b)~$b/a = 1 - 7/10 = 0.3$。(c)~分类时
采用视轴比：扁平椭球体沿短轴方向看去
轮廓呈圆形，故归为 E0。
\end{solution}

\begin{solution}{exo:b3:milky-way-galaxies:4}
$d = 10^{\mu/5 + 1}$~pc：大麦哲伦云 \qty{49.6}{kpc}，仙女座
\qty{783}{kpc}，NGC~4258 \qty{7.58}{Mpc}。大麦哲伦云中的一颗恒星
视差为 $1/\qty{49.6}{kpc} = \qty{0.020}{mas}$，约为最好星表
对亮星的不确定度：单个视差根本
给不出距离，大麦哲伦云须借助其他方法（食
双星）确定，\omterm{def:b3:milky-way-galaxies:cepheid}{造父变星}的距离标度由此出发。
\end{solution}

\begin{solution}{exo:b3:milky-way-galaxies:5}
(a)~$\omega_0 = A - B = \qty{27.2}{km.s^{-1}.kpc^{-1}}$，而 $V_0 =
27.2\times8.178 = \qty{222}{km.s^{-1}}$. (b)~$V'_0 = -(A + B) =
\qty{-3.4}{km.s^{-1}.kpc^{-1}}$：曲线平缓下降。(c)~两者都
均低于红巨星的数值（229 与 \qty{-1.7}{km.s^{-1}} 每
千秒差距）。该曲线的作者也得到同样的结论：$V_0$ 约低两
个标准差（由 $R_0$ 与这些\omterm{def:b3:milky-way-galaxies:rotation-curve}{奥尔特常数}推得）。
奥尔特公式假设轴对称圆盘中为圆轨道；非零的 $C$ 与
$K$ 表明局地速度场受棒状结构与
旋臂扰动，局地的 $A$ 与 $B$ 也携带这些扰动。
\end{solution}

\begin{solution}{exo:b3:milky-way-galaxies:6}
(a)~$\rho = V^2/(4\pi GR^2) = 229^2/(4\pi\times\num{4.301e-6}\times
8.178^2) = \num{1.45e7}$ 太阳质量每立方千秒差距，或
每立方秒差距 \num{0.0145} 个太阳质量。(b)~$M_\odot c^2 = \num{1.9884e30}\times
(\num{2.998e8})^2 = \qty{1.787e47}{J} = \qty{1.115e57}{GeV}$，且
$\qty{1}{pc^3} = (\qty{3.0857e18}{cm})^3 = \qty{2.938e55}{cm^3}$：比值
为 \qty{37.96}{GeV.cm^{-3}}。于是 $0.30/37.96 = \num{0.0079}$ 太阳
质量每立方秒差距。(c)~平坦曲线对应的密度计入了维持轨道的全部质量，
其中包括恒星与气体，仿佛它均匀分布于
球体之中；暗物质仅占其中一部分，此处约占一半。
\end{solution}

\begin{solution}{exo:b3:milky-way-galaxies:7}
(a)~$W_I(\mathrm{LMC}) = -3.314\times1.4771 + 15.888 = \qty{10.993}{mag}$,
故 $M_W = 10.993 - 18.477 = \qty{-7.484}{mag}$。(b)~$\mu = 25.00 + 7.484 =
32.484$，而 $d = 10^{32.484/5 + 1}~\mathrm{pc} = \qty{31.4}{Mpc}$。
(c)~$\delta d/d = (\ln 10/5)\,\delta\mu = 0.4605\times0.077 =
\qty{3.5}{\percent}$ 对单颗恒星而言，$3.5/\sqrt{50} = \qty{0.5}{\percent}$
对五十颗而言，若它们的误差彼此独立（大麦哲伦云自身距离的误差
距离的误差对所有恒星同样叠加）。
\end{solution}

\begin{solution}{exo:b3:milky-way-galaxies:8}
$\mu = I - A_I - M_I = 26.00 - 0.05 + 4.05 = 30.00$，故 $d =
10^{7}~\mathrm{pc} = \qty{10.0}{Mpc}$。误差 \qty{0.04}{mag} 相当于
$0.4605\times0.04 = \qty{1.8}{\percent}$ 的距离误差。
\end{solution}

\begin{solution}{exo:b3:milky-way-galaxies:9}
(a)~$d = cz/H_0$ 给出 $5\log_{10}d = 5\log_{10}(cz) - 5\log_{10}H_0$，
故 $m_B - M_B - 25 = 5\log_{10}(cz) - 5\log_{10}H_0$；代入
$\log_{10}(cz) = a_B + 0.2\,m_B$ 与 $5\log_{10}H_0 = 5a_B + M_B + 25$，此式
即为所求结果。(b)~$\log_{10}H_0 = -3.8506 + 0.714158 + 5 = 1.8636$，故
$H_0 = \qty{73.04}{km.s^{-1}.Mpc^{-1}}$. (c)~$M_B = 5(\log_{10}67.4 -
5.714158) = \qty{-19.427}{mag}$：超新星偏亮 \qty{0.174}{mag}，
因而所有距离放大的因子为 $10^{0.2\times0.174} = 1.084$，
恰为比值 $73.04/67.4$。\omterm{def:b3:milky-way-galaxies:tension}{哈勃张力}
体现为距离标度上约百分之八的分歧。
\end{solution}

\begin{solution}{exo:b3:milky-way-galaxies:10}
(a)~位于边缘、做圆轨道运动的恒星满足 $V^2/R \approx GM/R^2$
（球对称近似）。(b)~$R = \sqrt{M/(\pi\Sigma)}$，故 $V^2 =
GM/R = G\sqrt{\pi\Sigma}\,M^{1/2}$，而 $M = V^4/(\pi\Sigma G^2) \propto
V^4$；若质光比相同，则 $L \propto V^4$ 亦然。
(c)~$(100/229)^4 = \num{0.036}$：速度较快的星系质量
是另一个的 27 倍。(d)~光度相同时，流量按 $1/d^2$ 下降，故较暗者
距离为另一个的两倍。
\end{solution}

\begin{solution}{exo:b3:milky-way-galaxies:11}
(a)~$\qty{1}{km.s^{-1}} = \qty{1.0227}{kpc.Ga^{-1}}$，故
\qty{229}{km.s^{-1}} 为 \qty{234.2}{kpc.Ga^{-1}}。(b)~$2\pi R/V$：
\qty{107}{Ma} 在 \qty{4}{kpc} 处，\qty{322}{Ma} 在 \qty{12}{kpc} 处。
(c)~$\Delta\phi = Vt\,(1/4 - 1/12) = 234.2\times0.1667 = \qty{39.0}{rad}$,
6.2 圈：旋臂将是在 4 与
\qty{12}{kpc} 之间紧绕六圈的螺线，远比观测所见的旋臂紧密，且每
多转一圈便更紧。(d)~在密度波中，图样以一个角速度转动，
如同刚体；恒星与气体穿过它，因而
旋臂的形状不会被卷绕起来。
\end{solution}

\begin{solution}{exo:b3:milky-way-galaxies:12}
(a)~$a = (4.2\times10^6\times16^2)^{1/3} = \qty{1024}{au}$；在
\qty{8178}{pc} 处，其张角为 $1024/8178 = \ang{;;0.125}$（一天文单位在一秒差距处
张角为一角秒）。(b)~$R_S = 2GM/c^2 = 2\times\num{6.674e-11}
\times\num{8.35e36}/(\num{2.998e8})^2 = \qty{1.24e10}{m} = \qty{0.083}{au}$;
直径 $2R_S$ 的张角为 $0.166/8178 = \num{2.0e-5}$ 角秒，约
\qty{20}{\micro as}。(c)~$R_S = \qty{4.02e-7}{pc}$，而 $M/(\tfrac43\pi
R_S^3) = \num{1.5e25}$ 太阳质量每立方秒差距，为
由人马座~A* 的运动与大小所设下界的四十倍：该界限
已与黑洞本身的密度只差四十倍以内。
\end{solution}

\begin{solution}{pb:b3:milky-way-galaxies:1}

\textbf{第一部分}

\textbf{1.} $v_r = A\,d\sin2\ell$：\ang{45} 处为 $+A d$，\ang{90} 处为 0，
\ang{135} 处为 $-A d$，\ang{180} 处为 0。

\textbf{2.} $A = (15.4 + 15.2)/(2\times\qty{1.00}{kpc}) =
\qty{15.3}{km.s^{-1}.kpc^{-1}}$.

\textbf{3.} 在 \ang{90} 处，$v_t = d(B - A) = 4.74\times(-5.91)\times1.00 =
\qty{-28.02}{km.s^{-1}}$，故 $B = 15.3 - 28.02 =
\qty{-12.72}{km.s^{-1}.kpc^{-1}}$.

\textbf{4.} 在 \ang{0} 处，$v_t = d(A + B) = 4.74\times0.55 =
\qty{2.61}{km.s^{-1}}$，故 $B = 2.61 - 15.3 = \num{-12.69}$：两者
在 \num{0.03} 以内相符。

\textbf{5.} $\omega_0 = A - B = \qty{28.0}{km.s^{-1}.kpc^{-1}}$，且
$V_0 = 28.0\times8.18 = \qty{229}{km.s^{-1}}$.

\textbf{6.} $V'_0 = -(A + B) = \qty{-2.6}{km.s^{-1}.kpc^{-1}}$：缓慢下降。

\textbf{7.} $P = 2\pi/\omega_0$: $2\pi\times\qty{8.18}{kpc}/\qty{229.2}{km.s^{-1}}
= \qty{219}{Ma}$；$4540/219 = 20.7$ 圈。

\textbf{第二部分}

\textbf{8.} 切点要求视线经过比太阳更靠近中心的位置，
即 $\sin\ell > 0$ 且 $\ell < \ang{90}$（或
镜像象限，$\ang{270} < \ell < \ang{360}$，此时最大
速度最为负值）。在太阳圈之外，\omterm{def:b3:remote-sensing-gis:insar}{视线方向}上到中心的距离
处处增大，因此没有哪一点被单独指出，
且仅凭速度无法判断气体的位置。

\textbf{9.} $R = 8.18\sin\ell$: \qtylist{5.26;6.27;7.08;7.69}{kpc}.

\textbf{10.} $V = v_\mathrm{max} + V_0\sin\ell$: $82.9 + 147.2 = 230.1$,
$55.6 + 175.4 = 231.0$, $32.4 + 198.3 = 230.7$, $14.9 + 215.2 =
\qty{230.1}{km.s^{-1}}$.

\textbf{11.} $\Delta\nu = \nu v/c = \num{1.420406e9}\times82.9/299792 =
\qty{393}{kHz}$，向低频方向偏移：气体正在远离。

\textbf{12.} 距红巨星约 \qty{3}{km.s^{-1}} 以内：介于 5 与
\qty{8}{kpc} 之间时曲线是平坦的。

\textbf{13.} 若为点质量，则 $V = V_0\sqrt{R_0/R} = 229\times\sqrt{8.18/5.26} =
\qty{286}{km.s^{-1}}$，在第一切点处，远快于
所求得的 \qty{230}{km.s^{-1}}：质量并不集中在中心，
大部分质量位于 5 至 \qty{8}{kpc} 之间及更远处。

\textbf{第三部分}

\textbf{14.} $\num{1.067e10} + 2\times\num{3.944e10} = \num{8.96e10}$ 太阳
质量，几乎是综述值 \num[separate-uncertainty=true]{5 \pm 1 e10} 的两倍；模型中的恒星
可能偏重，这会使下文求得的暗物质
被低估。

\textbf{15.} $V_b^2 = GMR^2/(R^2 + b^2)^{3/2} = \num{4.301e-6}\times\num{1.067e10}
\times400/(400.09)^{3/2}$, $V_b = \qty{47.9}{km.s^{-1}}$.

\textbf{16.} 薄盘：$(a + b)^2 = 30.80$，$V^2 = \num{4.301e-6}\times\num{3.944e10}
\times400/(430.8)^{3/2}$，$V = \qty{87.1}{km.s^{-1}}$。\omterm{def:b3:milky-way-galaxies:structure}{厚盘}：
$(a + b)^2 = 11.56$, $V = \qty{90.1}{km.s^{-1}}$.

\textbf{17.} $V_\mathrm{vis} = \sqrt{47.9^2 + 87.1^2 + 90.1^2} =
\qty{134.2}{km.s^{-1}}$.

\textbf{18.} 在 \qty{8.18}{kpc} 处：74.8、108.4 与 \qty{127.8}{km.s^{-1}}，
合计 \qty{183.5}{km.s^{-1}}，与 $V_0 = \qty{229}{km.s^{-1}}$ 相比：
可见物质在太阳处已然不足。

\textbf{19.} $\sqrt{G\times\num{8.96e10}/20} = \qty{138.8}{km.s^{-1}}$。
模型中的 \qty{134.2}{km.s^{-1}} 与之接近，这是因为在 \qty{20}{kpc} 处，
约位于四至六个盘标长度之外，几乎所有可见质量都位于
轨道之内：可见物质的曲线已变为开普勒式。

\textbf{第四部分}

\textbf{20.} $M = 20\times200^2/\num{4.301e-6} = \num{1.860e11}$ 太阳质量。

\textbf{21.} $M_\mathrm{vis} = 20\times134.2^2/\num{4.301e-6} = \num{8.37e10}$
太阳质量。

\textbf{22.} $M_\mathrm{dark} = \num{1.860e11} - \num{8.37e10} = \num{1.02e11}$;
两质量均为 $RV^2/G$，且 $R$ 相同，故所占比例为
$1 - (V_\mathrm{vis}/V)^2 = 1 - (134.2/200)^2 = \num{0.550}$.

\textbf{23.} $x = 20/14.8 = 1.351$, $\ln(2.351) - 1.351/2.351 = 0.8549 - 0.5748
= 0.2801$, $M = 4\pi\times\num{1.06e7}\times14.8^3\times0.2801 =
\num{1.21e11}$ 太阳质量；暗物质所占比例为
$1.21/(1.21 + 0.837) = \num{0.59}$。\omterm{def:b3:milky-way-galaxies:halo}{暗物质晕}拟合的是整条
曲线，但在 \qty{20}{kpc} 处模型偏高：该处的测量值
比光滑模型（\qty{209}{km.s^{-1}}）低约 \qty{10}{km.s^{-1}}。

\textbf{24.} $x = 8.18/14.8 = 0.553$, $\rho = \num{1.06e7}/(0.553\times1.553^2)
= \num{7.95e6}$ 太阳质量每立方千秒差距，合 \num{0.0080} 每
立方秒差距，或 \qty{0.30}{GeV.cm^{-3}}，即拟合作者引用的
数值。

\textbf{25.} $1 - 8.37/18.60 = \num{0.550}$：在
\qty{20}{kpc} 以内的质量中，暗物质所占比例为
\textbf{\qty[text-series-to-math=true]{55}{\percent}}.
\end{solution}
